% !TEX root = ../Thesis.tex
\section{Parameters for Overview of Verification Tools}
\label{parameters}

This chapter introduces the parameters, on which the tools get analyzed. These parameters shall provide the reader with a brief overview of the tool and its key features and properties. They also shall enable the reader to be able to compare these tools against each other.

The following parameters of the tools are analyzed:

\textbf{Usability}

Does the tool work fully autonomous or does it need some manual steps? How to integrate Rust source code into the tool?

\textbf{Technology}

What technology does the tool for generating and validating the proofs? How powerful is this technology? Can it be customized to specific needs, e.g. a special SMT solver?

\textbf{Feedback}

How does the tool report errors or success? If so, does the tool return evidence to the contrary?

\textbf{Language support and verification conditions}

Is the tool supporting all language features of Rust? What are its limitations? What kind of verification conditions does the tool support?

\textbf{Development}

Who has developed the tool? Is it backed by a company?

\textbf{Activity}

Is the tool in active development or abandoned? Is it feature complete? Does the tool get used in production?

\textbf{Scientific activity}

Does the tool often get quoted in scientific publications, workshops or presentations?

\textbf{Documentation}

Is there a viable documentation available for the tool?